\section{物理验证与可信依据}

数值复现的风险不在「代码能跑」，而在「结果凑巧像对的」。可信的依据有三条，相互独立，任何一条不过都不放过。

\subsection{依据一：多近似方法对比}

同一物理量若被两条独立路径算出相同结果，就排除「单一路径的 bug 碰巧出对」。

\begin{itemize}
  \item \textbf{球体：Table 2 体积分 vs Mie 级数}。Table 2 对电流做体积分，Mie 是球谐级数求和，数学路径完全不同，四通道在 $<0.21\%$（介电球）/ $<0.03\%$（金球）内一致。这直接锁定 Table 2 公式与实现正确。
  \item \textbf{Grahn：双路径 + 独立库}。Path A（长波）与 Path B（有限核）分别对 Table 1/Mie，误差 $5.7\times10^{-7}$/$2.9\times10^{-4}$；再对 \texttt{miepython}（$1.6\times10^{-15}$）与独立远场投影（$7.7\times10^{-15}$）。四个独立来源互证。
  \item \textbf{Fig.~3：有限元 vs 边界元}。同一几何、同一材料，两种数值方法求解：频域有限元（900G）与边界元（128G）。PEC 近似下 ED/MQ 在 $73\%/95\%$ 内一致，验证两套求解器与多极矩后处理；MD/EQ 的差异源于 PEC 丢损耗，是已知模型近似，非求解错误。
\end{itemize}

\subsection{依据二：极限退化}

精确公式在极限下必须退化回更简单的已知结果，这是对公式适用范围的强检验。

\begin{itemize}
  \item \textbf{Table 2 $\to$ Table 1}：$kr\to0$ 时 $j_l(kr)\to(kr)^l/(2l+1)!!$，精确多极矩逐条退化回长波近似，ED/MD/EQ/MQ 四条相对差 $<1\%$。这是推导最核心的退化检查。
  \item \textbf{Mie $\to$ Rayleigh}：小尺寸散射截面 $\propto x^4$，log--log 斜率 $4.00$（目标 $4\pm0.05$）。
  \item \textbf{大尺寸 $Q_{\rm ext}\to2$}：$\|Q_{\rm ext}(120)-2\|<0.05$。
  \item \textbf{Grahn Rayleigh slope}：$6.098$，落在 $[6.0,6.3]$。
  \item \textbf{Fig.~3 双盘}：无解析长波目标（非球形），退化未做，如实标注。
\end{itemize}

\subsection{依据三：能量守恒与光学定理}

这两条是 Maxwell 方程的直接推论，任何正确解都必须满足，是最「硬」的物理约束。

\begin{itemize}
  \item \textbf{能量守恒 $C_{\rm ext}=C_{\rm sca}+C_{\rm abs}$}：介电球无吸收，$C_{\rm ext}=C_{\rm sca}$ 在多 $x$ 扫描下误差 $<10^{-10}$。
  \item \textbf{光学定理 $S(0)$}：双路独立实现（级数求和 / 前向极限），复振幅差 $1.2\times10^{-14}$；金球 550/1000/1700 nm 吸收支路 $C_{\rm abs}>0$ 正确。
  \item \textbf{Grahn Eq.~(22) vs Eq.~(20)}：消光与散射在 $x=0.5$ 相对差 $5.96\times10^{-7}$。
  \item \textbf{Fig.~3 功率闭合}：$C_{\rm sca,flux}$ 与多极矩之和闭合，power balance $=1.0$，未解析比例 $4.7\times10^{-5}$（0.5 网格）。
\end{itemize}

\begin{important}
  三条依据的分工：多近似对比排除实现 bug；极限退化排除公式抄错；能量守恒/光学定理排除物理不自洽。球体与 Grahn 三条全过；Fig.~3 能量守恒与多近似（FEM/BEM）通过，极限退化因无解析目标待补。
\end{important}

完整数字与证据位置见附录 B。
